\chapter{Etichetta monastica}
\label{Etichetta}

\section{La riverenza e le posture}

Inchianarsi, o prostrarsi, davanti a ciò che si considera degno del massimo rispetto, che si tratti di anziani, maestri, oggetti o luoghi sacri, è un'usanza millenaria della cultura buddista e una delle tante belle forme di etichetta buddista.

Eseguita in modo corretto e consapevole, la prostrazione non è solo un rituale vuoto che implica l'umiliarsi davanti a immagini religiose per una cieca accettazione della consuetudine. Piuttosto, è un modo per mostrare fede, rispetto e gratitudine verso ciò che è veramente degno di tale venerazione. Inoltre, è un atto di umiltà, una qualità che è sia una benedizione per gli altri, sia un modo per inclinare la propria mente verso gli ideali e le aspirazioni spirituali che quegli oggetti rappresentano. Inchinarsi al Triplice Gioiello significa elevare il proprio cuore e la propria mente.

La forma e il metodo corretto per eseguire la prostrazione sono i seguenti:

\subsection{La riverenza o prostrazione a cinque punti}

Il termine ‘prostrazione a cinque punti’ si riferisce ai cinque punti del corpo che entrano in contatto con il suolo quando si esegue una prostrazione. Essi sono: la fronte, entrambi i gomiti e entrambe le ginocchia.

\captionsetup{font=small, labelformat=empty, position=above}

\subsection{La postura preparatoria per la riverenza:}

\begin{figure}[h]
	\centering
	
	\begin{minipage}{0.40\textwidth}
		\centering
		\caption{Gli uomini si inginocchiano con i palmi delle mani appoggiate sulle cosce.}
		\includegraphics[scale=0.25,width=1\textwidth]{prep-man-ai.jpg}
		
	\end{minipage}
	\hfill
	\begin{minipage}{0.40\textwidth}
		\centering
		\caption{Le donne adottano la stessa postura, ma con il collo del piede rivolto verso il pavimento.}
		\includegraphics[scale=0.25,width=1\textwidth]{prep-woman-ai.jpg}
		
	\end{minipage}
	
\end{figure}


\subsection{L' ‘Añjali’}

Giungere i palmi delle mani all'altezza del petto, rivolti verso l'alto con un'angolazione di quarantacinque gradi e con le dita unite.
\begin{figure}[h]
	\centering
	
	\begin{minipage}{0.40\textwidth}
		\centering
		\includegraphics[scale=0.25,width=1\textwidth]{anjali-man-ai.jpg}
		%\label{fig:1}
	\end{minipage}
	\hfill
	\begin{minipage}{0.40\textwidth}
		\centering
		\includegraphics[scale=0.25,width=1\textwidth]{anjali-woman-ai.jpg}
		%\label{fig:2}
	\end{minipage}
	
\end{figure}

\subsection{Il ‘Vandana’}

Portare l'añjali (i palmi uniti) alla fronte, in modo che i pollici si trovino tra le sopracciglia e le punte degli indici siano all'attaccatura dei capelli.

\begin{figure}[h]
	\centering
	
	\begin{minipage}{0.40\textwidth}
		\centering
		\caption{Gli uomini mantengono la testa dritta.}
		\includegraphics[scale=0.25,width=1\textwidth]{vandana-man-ai.jpg}
		%\label{fig:1}
	\end{minipage}
	\hfill
	\begin{minipage}{0.40\textwidth}
		\centering
		\caption{Le donne inclinano la testa leggermente in avanti.}
		\includegraphics[scale=0.25,width=1\textwidth]{vandana-woman-ai.jpg}

		%\label{fig:2}
	\end{minipage}
\end{figure}
\enlargethispage{3\baselineskip}
\subsection{L' ‘Abhivada’}
Mantenendo la schiena dritta, inchinarsi portando la fronte a terra in modo che tocchi lo spazio tra i palmi (a una distanza di circa una spanna), che ora saranno rivolti verso il basso, appoggiati sul pavimento.
\begin{figure}[h]
	\centering
	\begin{minipage}{0.40\textwidth}
		\centering
		\includegraphics[scale=0.2,width=1\textwidth]{abhivada-man-ai.jpg}
		\caption{Gli uomini mantengono gomiti e ginocchia allineati e a contatto.}
		%\label{fig:1}
	\end{minipage}
	\hfill
	\begin{minipage}{0.40\textwidth}
		\centering
		\includegraphics[scale=0.2,width=1\textwidth]{abhivada-woman-ai.jpg}
		\caption{Le donne posizionano i gomiti ai lati delle ginocchia.}
		%\label{fig:2}
	\end{minipage}
\end{figure}

\subsection{\thaiFont พับเพียบ ‘papiyap’}

\thaiFont{‘พับเพียบ’} \normalfont(pronunciato ‘papiyap’) è un termine thailandese che si riferisce alla postura tradizionale da adottare durante il canto. Sedersi in ‘papiyap’ significa sedersi con le gambe piegate su un lato, come mostrato nella figura sottostante. La posizione a gambe incrociate o del ‘loto’ è riservata alla pratica formale della meditazione. Sedersi con le gambe o i piedi rivolti verso un'immagine del Buddha o verso i monaci è considerato irrispettoso, poiché i piedi sono la parte più bassa del corpo. Pertanto, ‘papiyap’ è una postura appropriata da assumere durante il canto. Tuttavia, quando si cantano i versi di omaggio al Triplice Gioiello (Sezione 1-Vol. 1), ci si inginocchia sui talloni.
\enlargethispage{3\baselineskip}
\begin{figure}[h]
	\centering
	
	\begin{minipage}{0.40\textwidth}
		\centering
		\includegraphics[scale=0.25,width=1\textwidth]{papiap-man-ai.jpg}
		\caption{Gli uomini si siedono con una gamba piegata di lato e l'altra incrociata davanti, con la pianta del piede che tocca il ginocchio.}
		%\label{fig:1}
	\end{minipage}
	\hfill
	\begin{minipage}{0.40\textwidth}
		\centering
		\includegraphics[scale=0.25,width=1\textwidth]{papiap-woman-ai.jpg}
		\caption{Le donne si siedono con una gamba piegata di lato, ma con l'altra gamba piegata sotto di sé, mantenendo le ginocchia unite.}
		%\label{fig:2}
	\end{minipage}
	
\end{figure}

\subsection{Cambio de postura}


Non è sbagliato cambiare posizione quando il dolore si fa acuto, ma bisogna farlo in modo educato. Il modo migliore per cambiare posizione, sia seduti in meditazione che durante il canto, è spostare la posizione delle gambe da dietro, piegandosi in avanti e usando una mano a terra come appoggio.
\clearpage

\section{Etiqueta general en el monasterio y la sala de meditación}

Il monastero e, in particolare, la sala di meditazione principale sono uno spazio sacro. La sala di meditazione (sala Uposatha) deve sempre fornire, sia alla comunità monastica che a quella laica, un ambiente adeguato e propizio per la pratica della virtù, della meditazione e dello sviluppo della saggezza. Esistono alcuni tipi di comportamento che contribuiscono a creare un'atmosfera di buona pratica spirituale e di tranquilla riflessione, e altri tipi di comportamento che possono distruggerla rapidamente. È quindi importante che chiunque entri nel monastero e, per estensione, nella sala di meditazione, sia consapevole di come il proprio comportamento possa influenzare coloro che lo circondano. Ciò andrà a beneficio sia della propria pratica di Sati, sia dell'ambiente del monastero nel suo complesso.

Quando si entra nella sala principale, si prega di osservare le seguenti pratiche:


\begin{enumerate}
\item Se è necessario parlare con qualcuno, è meglio uscire per farlo. O almeno farlo nel modo più silenzioso e breve possibile, per non disturbare le persone intorno che potrebbero essere intente a meditare.

\item Non sono ammessi cappelli, cappucci, berretti, scarpe, cibo o bevande (comprese le bottiglie d'acqua) nella sala.

\item Si prega i genitori con neonati o bambini piccoli di portare i propri figli fuori dalla sala se iniziano a piangere o a comportarsi in modo incontrollabile durante le sessioni di meditazione e i discorsi di Dhamma.

\item Foto e video possono essere realizzati solo con il permesso diretto del monaco più anziano.

\item Il monastero è un luogo di celibato. La separazione dei sessi è praticata rigorosamente durante i ritiri di meditazione e qualsiasi forma di contatto fisico inappropriato con il sesso opposto deve essere evitata in qualsiasi punto del terreno del monastero.

\item Entrando nel monastero, le persone dovrebbero evitare qualsiasi tipo di abbigliamento che sia rivelatore o impudico.

\item Si prega di non sedersi puntando i piedi verso l'altare o verso i monaci, poiché ciò è considerato una mancanza di rispetto.

\item Tutti i telefoni cellulari o qualsiasi altro dispositivo elettrico che possa fare rumore devono essere lasciati fuori dalla sala o spenti prima di entrarvi.

\item Queste regole si applicano non solo alla sala di meditazione principale, ma a qualsiasi edificio del monastero in cui si tengono discorsi di Dhamma o si pratica la meditazione.
\end{enumerate}
Un atteggiamento di base di moderazione e rispetto è tutto ciò che serve.

Grazie.
